\item Un oscilador armónico simple de amplitud $A$ tiene una energía total $E$. Determine:
\begin{enumerate}
    \item la energía cinética y
    \item la energía potencial cuando la posición es un tercio de la amplitud.
    \item ¿Para qué valores de la posición la energía cinética es igual a la mitad de la energía potencial?
    \item ¿Hay algún valor de la posición donde la energía cinética es mayor que la máxima energía potencial? Explique.
\end{enumerate}